% !TEX root = ../../main.tex
% 由 chapter2.tex 输入。


\subsection*{元素的阶、循环性与置换生成（习题1--14）}

\begin{problem}
\textbf{习题1：阶互素时落入正规子群。}有限群$G$中，$N\trianglelefteq G$且$\gcd(|N|,|G/N|)=1$。若$a$的阶整除$|N|$，则$a\in N$。
\end{problem}
\begin{proof}[补充证明]
令$d=\operatorname{ord}(aN)$。由于$a^{\operatorname{ord}(a)}=e$，商群中也有$(aN)^{\operatorname{ord}(a)}=N$，故
\[
d\mid\operatorname{ord}(a)\mid|N|.
\]
另一方面，拉格朗日定理给出$d\mid|G/N|$。
两数互素，所以$d=1$，即$aN=N$，从而$a\in N$。
\end{proof}

\begin{problem}
\textbf{习题2：互素阶的分解。}若$\operatorname{ord}(c)=rs$、$\gcd(r,s)=1$，则$c=ab$，其中$a,b\in\langle c\rangle$的阶分别为$r,s$。
\end{problem}
\begin{proof}[补充证明]
由贝祖等式，取整数$u,v$使
\[
us+vr=1.
\]
令$a=c^{us}$、$b=c^{vr}$，则两者都是$c$的幂，且
\[
ab=c^{us+vr}=c.
\]
因为$us\equiv1\pmod r$，有$\gcd(u,r)=1$；同理$\gcd(v,s)=1$。
用幂的阶公式，
\begin{align*}
\operatorname{ord}(a)&=\frac{rs}{\gcd(rs,us)}=\frac{rs}{s}=r,\\
\operatorname{ord}(b)&=\frac{rs}{\gcd(rs,vr)}=\frac{rs}{r}=s.
\end{align*}
\end{proof}

\begin{problem}
\textbf{习题3：互素指数的开方。}若$\operatorname{ord}(a)=m$且$\gcd(k,m)=1$，则$x^k=a$在$\langle a\rangle$内有唯一解。
\end{problem}
\begin{proof}[补充证明]
由$\gcd(k,m)=1$，存在整数$u,v$使$ku+mv=1$。
取$x=a^u$，则
\[
x^k=a^{ku}=a^{1-mv}=a,
\]
故解存在。

若另有$y\in\langle a\rangle$满足$y^k=a$，写$y=a^t$，则
\[
kt\equiv1\pmod m.
\]
因$k$在模$m$下可逆，$t\equiv u\pmod m$，所以$y=a^u=x$，解唯一。
\end{proof}

\begin{problem}
\textbf{习题4：$ab$与$ba$同阶。}
\end{problem}
\begin{proof}[补充证明]
直接计算共轭：
\[
a^{-1}(ab)a=ba.
\]
因此对任意正整数$m$，
\[
(ba)^m=a^{-1}(ab)^m a,
\]
所以$(ab)^m=e$当且仅当$(ba)^m=e$。
两者若有有限阶，最小正指数相同；若没有这样的正指数，则均为无限阶。
\end{proof}

\begin{problem}
\textbf{习题5：给定阶元素的奇偶性。}有限群中，阶恰为$n>2$的元素个数为偶数。
\end{problem}
\begin{proof}[补充证明]
设$X=\{a\in G:\operatorname{ord}(a)=n\}$。
取逆把$X$映到自身，因为$a$与$a^{-1}$同阶。
若$a=a^{-1}$，则$a^2=e$，其阶整除2，与$n>2$矛盾。

因此取逆使$X$分为互不重叠的二元集合
\[
\{a,a^{-1}\}.
\]
$G$有限，故$X$有限，元素个数为偶数。
\end{proof}

\begin{problem}
\textbf{习题6：对称群的中心。}$Z(S_n)=\{e\}$，$n>2$。
\end{problem}
\begin{proof}[补充证明]
若$\sigma\in Z(S_n)$，它与每个对换$(ij)$交换，故
\[
(\sigma(i)\ \sigma(j))=\sigma(ij)\sigma^{-1}=(ij).
\]
因此$\sigma$保持每个二元子集$\{i,j\}$。

固定任意$i$。因为$n>2$，可以取互不相同的$i,j,k$，于是
\[
\sigma(i)\in\{i,j\}\cap\{i,k\}=\{i\}.
\]
所以$\sigma(i)=i$对所有$i$成立，即$\sigma=e$。
\end{proof}

\begin{problem}
\textbf{习题7：有理数加法群的有限生成子群。}$(\Q,+)$的任一有限生成子群都为循环群。
\end{problem}
\begin{proof}[补充证明]
设$H=\langle q_1,\ldots,q_t\rangle$。若所有$q_i=0$，则$H=\{0\}$，已为循环群。
否则取正公分母$D$，写
\[
q_i=\frac{m_i}{D},\qquad m_i\in\Z.
\]
令$d=\gcd(m_1,\ldots,m_t)>0$。任意整数线性组合满足
\[
\sum_{i=1}^t k_iq_i=\frac{\sum_i k_im_i}{D}\in\frac dD\Z,
\]
所以$H\subseteq(d/D)\Z$。

由贝祖等式可选整数$u_i$使$\sum_i u_im_i=d$，因此
\[
\frac dD=\sum_i u_iq_i\in H.
\]
故$(d/D)\Z\subseteq H$，最终$H=(d/D)\Z$循环。
这里的$\Q^+$表示有理数加法群，不表示正有理数集合。
\end{proof}

\begin{problem}
\textbf{习题8：有限阶生成元与交换性推出有限性。}交换群$G$由有限多个有限阶元素生成，则$G$有限。
\end{problem}
\begin{proof}[补充证明]
设$G=\langle a_1,\ldots,a_t\rangle$，并令$n_i=\operatorname{ord}(a_i)<\infty$。
交换性允许重排任意生成元乘积，所以每个元素都可写成
\[
a_1^{k_1}\cdots a_t^{k_t}.
\]
由$a_i^{n_i}=e$，可把指数分别模$n_i$化简，取
\[
0\leq k_i<n_i\qquad(1\leq i\leq t).
\]
这种表达未必唯一，但已给出元素数的上界
\[
|G|\leq\prod_{i=1}^t n_i<\infty.
\]
所以$G$有限。交换性正是能先把相同生成元集中起来的条件。
\end{proof}

\begin{problem}
\textbf{习题9$^*$：由幂集合刻画循环群。}令$G^k=\{a^k:a\in G\}$，$k\geq1$。
$G$为有限循环群，当且仅当它的每个子群都等于某个$G^k$。
\end{problem}
\begin{proof}[补充证明]
先说明有限性条件不能去掉。无限循环群$\langle a\rangle$中，任意正整数$k$均有$a^k\neq e$，故
\[
G^k\neq\{e\}.
\]
所以正整数幂集合不能表示它的平凡子群。下面证明所述的有限循环群版本。

\textbf{由循环性推出幂集合条件。}
若$G=\langle a\rangle$有限，则
\[
G^k=\{a^{kt}:t\in\Z\}=\langle a^k\rangle.
\]
每个子群都有这种形式；平凡子群可取$k=|G|$。

\textbf{由幂集合条件推出有限性。}
对平凡子群，可取正整数$E$使$G^E=\{e\}$。
因此$g^E=e$对所有$g$成立，各元素阶整除$E$，并且
\[
G^{k+E}=G^k.
\]
所以只有有限多个不同的幂集合，也只有有限多个子群。
每个循环子群有限；全群是这些有限多个有限循环子群的并，故$G$有限。

\textbf{证明每个西罗子群循环且唯一。}
对$p\mid|G|$，选择阶为最大$p$幂$p^r$的元素$a$。这样的元素存在，因为西罗第一定理给出$p$阶子群，且$G$有限。
由题设，$\langle a\rangle=G^k$对某个正整数$k$成立。
又$a\in G^k$，故$a=b^k$对某个$b\in G$成立。
若$p\mid k$，幂的阶公式
\[
\operatorname{ord}(b^k)
=\frac{\operatorname{ord}(b)}{\gcd(\operatorname{ord}(b),k)}
\]
说明$b$的阶中$p$的指数大于$r$；取$b$的适当幂即可得到更大$p$幂阶的元素，矛盾。
因此$p\nmid k$。

任取$p$幂阶元素$x$，设其阶为$p^s$。
取整数$u$使$ku\equiv1\pmod{p^s}$，就有
\[
x=(x^u)^k\in G^k=\langle a\rangle.
\]
所以全部$p$幂阶元素都属于$\langle a\rangle$。
取包含$\langle a\rangle$的西罗$p$子群$P$；反向每个$P$元素都是$p$幂阶，故$P\leq\langle a\rangle$。
于是$P=\langle a\rangle$，且任何西罗$p$子群都等于它。

各西罗子群正规且阶互素，因而交平凡并两两交换，给出$G$的内直和。
每个因子循环，互素阶生成元的积的阶为各阶的乘积$|G|$，所以$G$循环。
\end{proof}

\begin{problem}
\textbf{习题10：两组对换生成$S_n$。}$S_n$由$(12),(13),\ldots,(1n)$生成，也由$(12),(23),\ldots,(n-1\ n)$生成。
\end{problem}
\begin{proof}[补充证明]
\textbf{第一组。}任意对换涉及1时已在生成集中；若$i,j\neq1$，则
\[
(ij)=(1i)(1j)(1i).
\]
所以第一组生成全部对换，进而生成$S_n$。

\textbf{第二组。}对$i<j$，令
\[
u=(i\ i+1)(i+1\ i+2)\cdots(j-2\ j-1).
\]
当$j=i+1$时取$u=e$。有$u(j-1)=i,u(j)=j$，故
\[
(ij)=u(j-1\ j)u^{-1}.
\]
右边全部由相邻对换组成，所以第二组也生成全部对换，从而生成$S_n$。
\end{proof}

\begin{problem}
\textbf{习题11：两个生成元生成$S_n$。}$S_n=\langle(12),(12\cdots n)\rangle$，$n\geq2$。
\end{problem}
\begin{proof}[补充证明]
令$c=(12\cdots n)$，对$1\leq i<n$，共轭公式给出
\[
c^{i-1}(12)c^{-(i-1)}=(i\ i+1).
\]
因此$\langle(12),c\rangle$包含全部相邻对换。
由习题10，它包含$S_n$的生成集，故等于$S_n$。
\end{proof}

\begin{lemma}\label{p2:triple-attach}
设有一列3-轮换，从第一个轮换涉及的三个点开始，后续每个轮换涉及的三点都与已出现的点有交，且所有轮换涉及的点合起来为$\{1,\ldots,n\}$。那么这些轮换生成$A_n$。
\end{lemma}
\begin{proof}[补充证明]
（此引理为补充，用于展开习题12。）先说明加入一个新点的情况。
设已有一组点$U$上的全部偶置换，并有一个轮换$(abz)$，其中$a,b\in U,z\notin U$。
把$(abz)$取逆可交换前两个位置；再用$U$内的偶置换共轭，就可得到所有$(uvz)$，$u,v\in U$且$u\neq v$。
当$|U|=3$时，用三点轮换及取逆即可得到六个有序点对；当$|U|\geq4$时，先任取把$(a,b)$映为$(u,v)$的置换，若为奇置换，就在余下两个点交换以调整奇偶。

任意$U\cup\{z\}$上的3-轮换，要么完全在$U$内，要么为$(uvz)$或其逆。
因此生成的群含全部3-轮换，等于$U\cup\{z\}$上的交错群。

若一个新轮换与$U$只交一个点，写为$(cde)$，其中$c\in U$、$d,e\notin U$。
选$u\in U\setminus\{c\}$，用$U$内的偶置换共轭，可得$(ude)$。
现在$(cde)$已生成三点集合$\{c,d,e\}$上的交错群，而$(ude)$与它交两点。
由刚证明的情形，先得到$\{c,d,e,u\}$上的全部偶置换。
因此它包含$(cud)$、$(cue)$；将它们与原有$U$上的偶置换结合，逐次加入$d,e$，便得到$U\cup\{d,e\}$上的全部偶置换。

从最初三个点出发，按题设次序逐个加入轮换，归纳即可覆盖全部点，并生成$A_n$。
\end{proof}

\begin{problem}
\textbf{习题12：两个生成元生成$A_n$。}当$n>2$为偶数时，$A_n=\langle(123),(12\cdots n)^2\rangle$；当$n$为奇数时，$A_n=\langle(123),(34\cdots n)\rangle$。
\end{problem}
\begin{proof}[补充证明]
两种情形中所列生成元都是偶置换，所以生成的群包含于$A_n$。

\textbf{$n$为偶数。}令$c=(12\cdots n)$，对$j=0,1,\ldots,n/2-1$，
\[
c^{2j}(123)c^{-2j}=(2j+1\ 2j+2\ 2j+3),
\]
其中数字按模$n$取1至$n$之间的代表。
这些轮换涉及的三元集合可以依次加入，均与已出现的点有交，且最终覆盖全部$n$个点。
由引理\ref{p2:triple-attach}，它们生成$A_n$。

\textbf{$n$为奇数。}当$n\geq5$，令$d=(34\cdots n)$。
它固定1、2，且其幂将3依次映到$3,4,\ldots,n$，所以共轭得到
\[
d^j(123)d^{-j}=(1\ 2\ d^j(3)).
\]
因而生成群含全部$(12k)$，$3\leq k\leq n$。
这些轮换每次加入一个新点，由同一引理也生成$A_n$。
$n=3$时，$(34\cdots n)$按固定全部点的恒等置换处理，而$(123)$已生成$A_3$。
\end{proof}

\begin{problem}
\textbf{习题13：长轮换的中心化子。}若$\sigma=(12\cdots n)$，则$C_{S_n}(\sigma)=\langle\sigma\rangle$，其共轭类大小为$(n-1)!$。
\end{problem}
\begin{proof}[补充证明]
若$\tau\sigma=\sigma\tau$，则对所有$j\geq0$，
\[
\tau(\sigma^j(1))=\sigma^j(\tau(1)).
\]
因为$1,\sigma(1),\ldots,\sigma^{n-1}(1)$遍历全部点，$\tau$完全由$\tau(1)$确定。
取唯一的$0\leq k<n$使$\tau(1)=\sigma^k(1)$，上式便说明$\tau=\sigma^k$。
反过来，$\sigma$的幂均与$\sigma$交换，所以
\[
C_{S_n}(\sigma)=\langle\sigma\rangle.
\]
共轭作用的稳定子为此中心化子，故类大小为
\[
[S_n:C_{S_n}(\sigma)]=\frac{n!}{n}=(n-1)!.
\]
\end{proof}

\begin{problem}
\textbf{习题14：$S_5$的共轭类及正规子群。}
\end{problem}
\begin{solution}
\textit{补充解答。}
同一轮换类型的置换互相共轭，故各类代表及大小为
\[
\begin{array}{c|rrrrrrr}
\text{代表}&e&(12)&(12)(34)&(123)&(123)(45)&(1234)&(12345)\\\hline
\text{类大小}&1&10&15&20&20&30&24
\end{array}
\]
例如3-轮换有$\binom53\cdot2=20$个，双对换有$5\cdot3=15$个，5-轮换有$4!=24$个。
一般地，若类型含$m_i$个$i$-轮换，则可先排列全部符号，再除去每个轮换的$i$种循环起点和等长轮换间的$m_i!$种次序，因此
\[
\text{类大小}=\frac{5!}{\prod_i i^{m_i}m_i!}.
\]

由这些类判断正规子群。正规子群含$e$且为共轭类的并，阶须整除120。
在所列类大小的子集和中，候选阶只有1、40、60、120。
40只能为$1+15+24$；60只能为$1+15+20+24$，其中20暂有两个类可选。
这两种候选都包含全部双对换，而
\[
((12)(34))((12)(35))=(354),
\]
乘积为3-轮换。因此40阶候选不对乘法封闭，60阶候选中的20也必须选3-轮换类。
剩下的60阶候选恰为所有偶置换，确为$A_5$。
因此正规子群只有$\{e\},A_5,S_5$。
\end{solution}

\subsection*{子群、陪集、直和及小阶群（习题15--30）}

\begin{problem}
\textbf{习题15：商群中子群的交。}若$N\trianglelefteq G$且$N\leq H_1\cap H_2$，则
$(H_1/N)\cap(H_2/N)=(H_1\cap H_2)/N$。
\end{problem}
\begin{proof}[补充证明]
由于$N\leq H_i$，对任意$g\in G$，
\[
gN\in H_i/N\iff g\in H_i.
\]
所以
\begin{align*}
gN\in(H_1/N)\cap(H_2/N)
&\iff g\in H_1\cap H_2\\
&\iff gN\in(H_1\cap H_2)/N.
\end{align*}
两边在$G/N$内实际相等，故当然同构。
\end{proof}

\begin{problem}
\textbf{习题16：交子群的陪集。}$H_1\cap H_2$的每个左陪集可写成$H_1$一个左陪集与$H_2$一个左陪集的交。
\end{problem}
\begin{proof}[补充证明]
任取左陪集$a(H_1\cap H_2)$，左乘$a$为双射，故
\[
a(H_1\cap H_2)=aH_1\cap aH_2.
\]
具体地，左边元素显然在右边；若$x=ah_1=ah_2$，由消去律$h_1=h_2\in H_1\cap H_2$，所以右边也包含于左边。

更一般地，若$aH_1\cap bH_2$非空，取$x$在交中，则
\[
aH_1=xH_1,\qquad bH_2=xH_2,
\]
所以该交等于$x(H_1\cap H_2)$。
\end{proof}

\begin{problem}
\textbf{习题17：有限指数子群的交。}若$[G:H_1],[G:H_2]$有限，则
$[G:H_1\cap H_2]\leq[G:H_1][G:H_2]$。
\end{problem}
\begin{proof}[补充证明]
定义陪集集合之间的映射
\[
\theta:\mathcal L_G(H_1\cap H_2)\to\mathcal L_G(H_1)\times\mathcal L_G(H_2),
\qquad g(H_1\cap H_2)\longmapsto(gH_1,gH_2).
\]
这里的$\mathcal L_G(H)$均为左陪集集合，不要求$H$正规。
若两个源陪集相同，则$g^{-1}g'\in H_1\cap H_2$，两个目标陪集分别相同，故良定义。
若目标相同，反过来也有$g^{-1}g'\in H_1\cap H_2$，所以源陪集相同，映射单射。

右侧集合有限，计数得到
\[
[G:H_1\cap H_2]\leq[G:H_1][G:H_2]<\infty.
\]
\end{proof}

\begin{problem}
\textbf{习题18：有限指数给出正规子群或忠实表示。}有限群$G$有真子群$H$、$[G:H]=n>1$，则或有非平凡正规子群，其指数整除$n!$，或$G$同构于$S_n$的子群。
\end{problem}
\begin{proof}[补充证明]
令$G$作用于$H$的$n$个左陪集，得到同态$\rho:G\to S_n$，记$K=\ker\rho$。
若$g\in K$，它特别固定陪集$H$，所以$gH=H$，即$g\in H$。故
\[
K\trianglelefteq G,\qquad K\leq H<G.
\]
若$K\neq\{e\}$，它是非平凡真正规子群，且第一同构定理给出
\[
[G:K]=|\rho(G)|\mid|S_n|=n!.
\]
若$K=\{e\}$，则$\rho$单射，$G\cong\rho(G)\leq S_n$。
\end{proof}

\begin{problem}
\textbf{习题19：最小素因子指数的子群正规。}若$p$为$|G|$的最小素因子，$[G:H]=p$，则$H\trianglelefteq G$。
\end{problem}
\begin{proof}[补充证明]
让$G$作用于$H$的$p$个左陪集，记核为$K$，像为$Q\leq S_p$。
由上一题，$K\leq H$，并且
\[
|Q|\mid|G|,\qquad |Q|\mid p!.
\]
$|Q|$的素因子不得小于$p$，因为$p$是$|G|$最小素因子；但$p!$中唯一不小于$p$的素因子就是$p$，其指数为1。
所以$|Q|$只能为1或$p$。

陪集作用传递，其某个轨道大小为$p$，故$p\mid|Q|$，于是$|Q|=p$。
因此
\[
[G:K]=p=[G:H].
\]
结合$K\leq H$得到$K=H$，所以$H\trianglelefteq G$。
\end{proof}

\begin{problem}
\textbf{习题20：$p^2$阶群的分类。}这样的群交换，且恰有$C_{p^2}$、$C_p\times C_p$两种同构类型。
\end{problem}
\begin{proof}[补充证明]
\textbf{先证交换。}类方程给出$Z(G)\neq\{e\}$，其阶只能为$p$或$p^2$。
若为$p^2$，已经交换；若为$p$，则$G/Z(G)$阶为$p$，所以循环。
取生成元的提升$a$，任意元素都可写为$a^iz$、$z\in Z(G)$，因而
\[
(a^iz)(a^jw)=a^{i+j}zw=a^{j+i}wz=(a^jw)(a^iz).
\]
故这种情形$G$也交换。

\textbf{再分类。}若有$p^2$阶元素，群即为$C_{p^2}$。
否则所有非单位元都为$p$阶。取$a\neq e$，再取$b\notin\langle a\rangle$。
两个$p$阶循环子群交平凡，所以
\[
|\langle a\rangle\langle b\rangle|=p^2=|G|.
\]
它们交换，故乘法映射$C_p\times C_p\to G$为同构。
两种群由是否存在$p^2$阶元素区分，确实不同构。
\end{proof}

\begin{problem}
\textbf{习题21：非交换6阶群。}每个非交换6阶群同构于$S_3$。
\end{problem}
\begin{proof}[补充证明]
西罗计数给出$n_3\mid2,n_3\equiv1\pmod3$，所以$n_3=1$。
设唯一的3阶子群为$\langle a\rangle$，再由西罗第一定理取2阶元素$b$。
两个子群交平凡，其乘积大小为6，所以$G=\langle a,b\rangle$。

共轭$b$保持$\langle a\rangle$，只能把$a$映到$a$或$a^{-1}$。
若映到$a$，两个生成元交换，$G$就交换，与题设矛盾。
因此
\[
a^3=b^2=e,\qquad bab=a^{-1}.
\]
这使乘积化为$a^i$或$a^ib$，$0\leq i<3$，与$D_3$的六个正规形有相同乘法。
将$a$映到$(123)$、$b$映到$(12)$给出$G\cong S_3$。
\end{proof}

\begin{problem}
\textbf{习题22：15阶群。}全部15阶群同构于$C_{15}$。
\end{problem}
\begin{proof}[补充证明]
由西罗数约束，
\begin{align*}
n_5\mid3,\quad n_5\equiv1\pmod5&\quad\Longrightarrow\quad n_5=1,\\
n_3\mid5,\quad n_3\equiv1\pmod3&\quad\Longrightarrow\quad n_3=1.
\end{align*}
两个西罗子群分别为$C_5,C_3$，且都正规。
它们的交平凡，来自不同因子的元素交换，其乘积有15个元素，故
\[
G\cong C_3\times C_5.
\]
两个生成元的积阶为$3\cdot5=15$，所以该直积为$C_{15}$。
\end{proof}

\begin{problem}
\textbf{习题23：10阶群。}恰有$C_{10}$与$D_5$两种同构类型。
\end{problem}
\begin{proof}[补充证明]
由$n_5\mid2$且$n_5\equiv1\pmod5$，5阶子群唯一，记为$\langle a\rangle$。
取2阶元素$b$，两个子群乘积有10个元素，所以$G=\langle a,b\rangle$。
共轭给出
\[
bab=a^t,\qquad\gcd(t,5)=1.
\]
因$b^2=e$，再共轭一次得到$a=a^{t^2}$，即$t^2\equiv1\pmod5$。
在模5下只有$t=1,-1$。

若$t=1$，生成元交换，$G\cong C_5\times C_2\cong C_{10}$。
若$t=-1$，关系与二面体群一致，正规形$a^i,a^ib$给出$G\cong D_5$。
两者一个交换、一个非交换，所以不同构。
\end{proof}

\begin{problem}
\textbf{习题24：不同素数乘积阶的群不单。}若$|G|=pq$、$p<q$，则$G$不是单群。
\end{problem}
\begin{proof}[补充证明]
交换两个素数的名字，可设$p<q$。
西罗$q$子群个数满足
\[
n_q\mid p,\qquad n_q\equiv1\pmod q.
\]
它只能为1或$p$，但$p<q$排除$p\equiv1\pmod q$，所以$n_q=1$。
于是唯一的$q$阶子群为非平凡真正规子群，$G$不单。
\end{proof}

\begin{problem}
\textbf{习题25：$p^2q$阶群有正规西罗子群。}其中$p,q$为不同素数。
\end{problem}
\begin{proof}[补充证明]
若$n_p=1$，已经得到正规西罗子群。
否则由$n_p\mid q$可知$n_p=q$，而$n_p\equiv1\pmod p$给出$q>p$。

若$n_q=1$同样已完成。反设$n_q\neq1$，由$n_q\mid p^2$，只能为$p$或$p^2$。
由于$p<q$，前者不可能模$q$余1，故$n_q=p^2$，即
\[
q\mid p^2-1=(p-1)(p+1).
\]
$q$为素数且大于$p$，因此只能$q=p+1$。
两者都为素数迫使$p=2,q=3$，故只剩12阶情形。

若此时$n_3=4$，四个3阶子群的非单位元两两不重复，共给出$4\cdot2=8$个3阶元。
余下只有四个元素，包含$e$。
任一4阶西罗2子群不能包含3阶元，故必须等于这四个元素的集合，因而唯一，即$n_2=1$。
这与$n_p=q=3$矛盾，所以至少有一个西罗子群正规。
\end{proof}

\begin{problem}
\textbf{习题26：两个非交换的有限例子。}分别构造所有元素满足$x^3=e$、及所有元素满足$x^4=e$的有限非交换群。
\end{problem}
\begin{solution}
\textit{补充解答。}
\textbf{全部元素立方为单位元的例子。}在$\F_3$上取
\[
H=\left\{\begin{pmatrix}1&a&c\\0&1&b\\0&0&1\end{pmatrix}:a,b,c\in\F_3\right\}.
\]
矩阵乘法及求逆保持这一形状，故它为27阶群。
每个元素写为$I+N$，其中$N^3=0$；在特征3中，
\[
(I+N)^3=I+3N+3N^2+N^3=I.
\]
但
\begin{align*}
(I+E_{12})(I+E_{23})&=I+E_{12}+E_{23}+E_{13},\\
(I+E_{23})(I+E_{12})&=I+E_{12}+E_{23},
\end{align*}
所以$H$非交换。

\textbf{全部元素四次方为单位元的例子。}取$D_4$。
其旋转的阶整除4，各反射的阶为2，因此所有元素满足$x^4=e$，而$D_4$非交换。
这分别给出两个有限非交换群。
\end{solution}

\begin{problem}
\textbf{习题27：双陪集计数。}有限群中，$H,K\leq G,a\in G$，有
\[
|HaK|=|H|[K:a^{-1}Ha\cap K].
\]
\end{problem}
\begin{proof}[补充证明]
双陪集可以分成右陪集$Hak$，其中$k\in K$。
对$k_1,k_2\in K$，
\begin{align*}
Hak_1=Hak_2
&\iff ak_1k_2^{-1}a^{-1}\in H\\
&\iff k_1k_2^{-1}\in a^{-1}Ha\cap K.
\end{align*}
所以这些不同右陪集的个数，就是子群$L=a^{-1}Ha\cap K$在$K$中的右陪集个数$[K:L]$。
不同右陪集不交，每个大小均为$|H|$，从而
\[
|HaK|=|H|[K:L].
\]
\end{proof}

\begin{problem}
\textbf{习题28：西罗2子群只有2阶时的指数2子群。}
\end{problem}
\begin{proof}[补充证明]
使用左正则表示$\rho:G\to S(G)$。
设$|G|=2m$，其中$m$为奇数，取西罗2子群的非单位元$a$，则$a^2=e$。
左乘$a$没有不动点，因为$ax=x$会推出$a=e$；其置换将$G$分为$m$对
\[
\{x,ax\}.
\]
因此$\rho(a)$为$m$个不相交对换的积，符号为$(-1)^m=-1$。
故复合同态
\[
G\xrightarrow{\rho}S(G)\xrightarrow{\operatorname{sgn}}\{1,-1\}
\]
满射，其核便是指数2的子群。
\end{proof}

\begin{problem}
\textbf{习题29：大子集的乘积。}有限群中$A,B\subseteq G$，若$|A|+|B|>|G|$，则$AB=G$。
\end{problem}
\begin{proof}[补充证明]
固定任意$g\in G$。集合$gB^{-1}$与$B$大小相同，所以
\[
|A|+|gB^{-1}|>|G|.
\]
因此它们不能不交。选$a\in A\cap gB^{-1}$，则存在$b\in B$使$a=gb^{-1}$，从而
\[
g=ab\in AB.
\]
$g$任意，故$G\subseteq AB$；反向包含显然，得$AB=G$。
\end{proof}

\begin{problem}
\textbf{习题30：真子群的共轭不能覆盖有限群。}若$H<G$，则$G\neq\bigcup_{g\in G}gHg^{-1}$。
\end{problem}
\begin{proof}[补充证明]
设$|G|=n,|H|=h,[G:H]=m>1$。
不同共轭子群数为
\[
r=[G:N_G(H)]\leq[G:H]=m.
\]
这些子群都含$e$，所以先只计一次$e$，其余元素的并至多有$r(h-1)$个。
因此
\begin{align*}
\left|\bigcup_{g\in G}gHg^{-1}\right|
&\leq1+r(h-1)\\
&\leq1+m(h-1)\\
&=n-m+1<n.
\end{align*}
所以此并集不能等于$G$。
\end{proof}

\subsection*{西罗子群、中心与可解性（习题31--45）}

\begin{problem}
\textbf{习题31：$S_4$的西罗2子群及西罗3子群。}
\end{problem}
\begin{solution}
\textit{补充解答。}
\textbf{西罗2子群。}$|S_4|=24=2^3\cdot3$，故其阶为8。
四个点分成两个无序二元组的方式恰有
\[
12|34,\qquad13|24,\qquad14|23.
\]
每个分块的稳定子有$2\cdot2\cdot2=8$个元素：可分别交换组内两点，也可交换两组。
这给出三个不同的8阶子群。
因为$n_2\mid3$，这三个已经穷尽全部西罗2子群。
例如$\langle(1234),(13)\rangle$稳定分块$13|24$，同构于$D_4$。

\textbf{西罗3子群。}其阶为3，每个由一个3-轮换生成。
共有$\binom43\cdot2=8$个3-轮换，每个3阶群含两个非单位元，所以共有四个：
\[
\langle(123)\rangle,\quad\langle(124)\rangle,
\quad\langle(134)\rangle,\quad\langle(234)\rangle.
\]
\end{solution}

\begin{problem}
\textbf{习题32：56阶与148阶群不单。}
\end{problem}
\begin{proof}[补充证明]
\textbf{148阶。}由$148=4\cdot37$，
\[
n_{37}\mid4,\qquad n_{37}\equiv1\pmod{37},
\]
所以$n_{37}=1$，唯一37阶子群为非平凡真正规子群。

\textbf{56阶。}由$56=8\cdot7$，$n_7$只能为1或8。
若为1已得结论。
若为8，八个7阶子群互交于$e$，共贡献
\[
8(7-1)=48
\]
个不同的7阶元。
因此不为7阶的元素只剩$56-48=8$个，包含$e$。
任意8阶西罗2子群的元素阶均为2的幂，故只能由这八个元素组成；每个西罗2子群都必须等于这个集合，因而唯一、正规。
两种情况下$G$均不单。
\end{proof}

\begin{problem}
\textbf{习题33：弗拉蒂尼论证。}有限群$G$中，$N\trianglelefteq G$，$P$是$N$的西罗$p$子群，则$G=N N_G(P)$。
\end{problem}
\begin{proof}[补充证明]
取任意$g\in G$。因为$N$正规，$gPg^{-1}$仍在$N$内，且与$P$同阶，故为$N$的西罗$p$子群。
在$N$内使用西罗共轭定理，可选$n\in N$使
\[
gPg^{-1}=nPn^{-1}.
\]
所以
\[
(n^{-1}g)P(n^{-1}g)^{-1}=P,
\]
即$n^{-1}g\in N_G(P)$。
于是$g=n(n^{-1}g)\in N N_G(P)$。
$g$任意，而反向包含显然，故$G=N N_G(P)$。
\end{proof}

\begin{problem}
\textbf{习题34：正规子群与商群的西罗子群。}若$P$是$G$的西罗$p$子群、$N\trianglelefteq G$，则$P\cap N$是$N$的西罗$p$子群，$PN/N$是$G/N$的西罗$p$子群。
\end{problem}
\begin{proof}[补充证明]
先取$N$的西罗$p$子群$Q$，并把它包含在$G$的某个西罗$p$子群$P_1$内。
再取$g\in G$使$gP_1g^{-1}=P$。由于$N\trianglelefteq G$，
\[
gQg^{-1}\leq P\cap N.
\]
左边已是$N$的西罗$p$子群，右边又是$N$的$p$子群，故阶必须相同，$P\cap N$为所求西罗子群。

写$|G|=p^rm$、$|N|=p^sn$，其中$p\nmid mn$。
前半已知$|P\cap N|=p^s$，由第二同构定理，
\[
|PN/N|=|P/(P\cap N)|=p^{r-s}.
\]
而$|G/N|$中的全部$p$因子恰为$p^{r-s}$，所以$PN/N$为商群的西罗$p$子群。
\end{proof}

\begin{problem}
\textbf{习题35：子群的西罗子群可由共轭后相交得到。}若$H\leq G$，$P$是$G$的西罗$p$子群，则存在$a\in G$使$aPa^{-1}\cap H$是$H$的西罗$p$子群。
\end{problem}
\begin{proof}[补充证明]
取$H$的西罗$p$子群$Q$。
由西罗第二定理，可选$a\in G$使
\[
Q\leq aPa^{-1}.
\]
因此$Q\leq aPa^{-1}\cap H$。
右边是$H$的$p$子群，而$Q$已为其西罗子群，所以不能有更大的$p$子群包含它。
故
\[
aPa^{-1}\cap H=Q,
\]
完成证明。
\end{proof}

\begin{problem}
\textbf{习题36：$\GL_n(\F_p)$的西罗$p$子群。}$n\geq1$。
\end{problem}
\begin{solution}
\textit{补充解答。}
取全体单位上三角矩阵组成的群
\[
U=\{(a_{ij}):a_{ii}=1,\ a_{ij}=0\ (i>j)\}.
\]
每个严格上三角位置可任取$p$个值，位置数为$n(n-1)/2$，所以
\[
|U|=p^{n(n-1)/2}.
\]
由前章可逆矩阵计数公式，
\begin{align*}
|\GL_n(\F_p)|
&=\prod_{k=0}^{n-1}(p^n-p^k)\\
&=p^{0+1+\cdots+(n-1)}\prod_{i=1}^n(p^i-1)\\
&=p^{n(n-1)/2}\prod_{i=1}^n(p^i-1).
\end{align*}
后一个乘积不被$p$整除，故$U$为西罗$p$子群。

下面证明它的正规化子为全部可逆上三角矩阵群$B$，从而
\[
n_p=[\GL_n(\F_p):B]
=\frac{\prod_{i=1}^n(p^i-1)}{(p-1)^n}.
\]
\end{solution}
\begin{proof}[补充证明]
上三角矩阵的共轭保持上三角形状，且对角线乘法逐分量进行，因此将对角线全1的矩阵仍映为此种矩阵。
故$B\leq N_G(U)$。

设$g$正规化$U$，先考察$U$的全部公共不动向量。
对$v=(v_1,\ldots,v_n)^{\mathsf T}$及$i<j$，
\[
(I+E_{ij})v=v+v_je_i.
\]
若$v$被全部$U$固定，依次取$i=1,j=2,\ldots,n$，得到$v_2=\cdots=v_n=0$。
反向$e_1$的所有倍数都被固定，所以公共不动向量集合为$L_1=\langle e_1\rangle$。

因为$gUg^{-1}=U$，$g$将公共不动向量仍映为公共不动向量，故保持$L_1$。
在商空间$\F_p^n/L_1$中，$U$诱导$(n-1)$维单位上三角群。
重复同一论证，得到$g$保持
\[
L_1\subset L_2\subset\cdots\subset L_n,
\qquad L_i=\langle e_1,\ldots,e_i\rangle.
\]
保持这些子空间恰好说明$g$为上三角矩阵，所以$N_G(U)=B$。

$B$的对角线有$(p-1)^n$种选择，严格上三角部分有$p^{n(n-1)/2}$种，故
\[
|B|=(p-1)^np^{n(n-1)/2}.
\]
代入西罗数的指数公式即得结论。$n=1$时$U=\{1\}$、$B=G$，公式给出$n_p=1$，仍成立。
\end{proof}

\begin{problem}
\textbf{习题37：非交换$p^3$阶群的换位子群。}若$|G|=p^3$且$G$非交换，则$G^{(1)}=Z(G)$，两者阶为$p$。
\end{problem}
\begin{proof}[补充证明]
$G$为非平凡$p$群，中心非平凡，其阶可能为$p,p^2,p^3$。
若阶为$p^3$，则$G$交换；若阶为$p^2$，则$G/Z(G)$阶为$p$而循环，由习题39的证明也推出交换。
这都与题设矛盾，所以
\[
|Z(G)|=p.
\]
商群$G/Z(G)$阶为$p^2$，由习题20交换，故$G^{(1)}\leq Z(G)$。
而$G$非交换使$G^{(1)}\neq\{e\}$；在$p$阶群中只有平凡及全体两个子群，因此
\[
G^{(1)}=Z(G).
\]
\end{proof}

\begin{problem}
\textbf{习题38：$p$群中的$p$阶正规子群在中心内。}若$N\trianglelefteq G$、$|N|=p$且$G$为有限$p$群，则$N\leq Z(G)$。
\end{problem}
\begin{proof}[补充证明]
正规性使共轭限制为$N$的自同构，从而得到同态
\[
\rho:G\to\Aut(N).
\]
$N$阶为$p$而循环，所以$|\Aut(N)|=p-1$。
$\rho(G)$的阶既整除$p$幂$|G|$，又整除$p-1$，所以为1。
因此对全部$g\in G,n\in N$，
\[
gng^{-1}=n,
\]
即$N\leq Z(G)$。
\end{proof}

\begin{problem}
\textbf{习题39：中心商群循环则原群交换。}
\end{problem}
\begin{proof}[补充证明]
设$G/Z(G)=\langle aZ(G)\rangle$。
任意$x,y\in G$都可写为
\[
x=a^mz,\qquad y=a^nw,\qquad z,w\in Z(G).
\]
于是
\begin{align*}
xy&=a^mza^nw\\
&=a^{m+n}zw\\
&=a^{n+m}wz=yx.
\end{align*}
故$G$交换。这里未使用有限性，结论对无限群也成立。
\end{proof}

\begin{problem}
\textbf{习题40：有限群嵌入某个交错群。}
\end{problem}
\begin{proof}[补充证明]
由凯莱定理，设$m=|G|$，有单射同态$G\hookrightarrow S_m$。
把$S_m$看成在$m+2$个点中固定末两点的置换群，令$t=(m+1\ m+2)$。
记$\varepsilon(\sigma)=0$或1，分别表示$\sigma$为偶置换或奇置换，定义
\[
\iota(\sigma)=\sigma t^{\varepsilon(\sigma)}.
\]
它总为偶置换。又$t$与$S_m$中元素交换，且$\varepsilon$在模2意义下相加，所以
\begin{align*}
\iota(\sigma)\iota(\tau)
&=\sigma\tau t^{\varepsilon(\sigma)+\varepsilon(\tau)}\\
&=\iota(\sigma\tau).
\end{align*}
若$\iota(\sigma)=e$，限制于前$m$个点即得$\sigma=e$，故$\iota$单射。
将两个嵌入复合，得到$G$同构于$A_{m+2}$的子群。
\end{proof}

\begin{problem}
\textbf{习题41：正规子群与换位子群交平凡时在中心内。}若$N\trianglelefteq G$且$N\cap G^{(1)}=\{e\}$，则$N\leq Z(G)$。
\end{problem}
\begin{proof}[补充证明]
对任意$n\in N,g\in G$，由$N$正规，有$g^{-1}ng\in N$，因此
\[
[n,g]=n^{-1}(g^{-1}ng)\in N.
\]
它同时是一个换位子，故也属于$G^{(1)}$。
由$N\cap G^{(1)}=\{e\}$，得到$[n,g]=e$，即$ng=gn$。
因$n,g$任意，$N$的每个元素都在中心，故$N\leq Z(G)$。
\end{proof}

\begin{problem}
\textbf{习题42：可解群的子群与商群仍可解。}
\end{problem}
\begin{proof}[补充证明]
设$G^{(r)}=\{e\}$。
对子群$H\leq G$，从$H^{(1)}\leq G^{(1)}$逐次取换位子群，得到
\[
H^{(r)}\leq G^{(r)}=\{e\}.
\]
所以$H$可解。

若$N\trianglelefteq G$，自然同态$\pi:G\to G/N$将换位子映为像的换位子；又每个像元素都有原像，故
\[
(G/N)^{(r)}=\pi(G^{(r)})=\{N\}.
\]
所以商群也可解。
\end{proof}

\begin{problem}
\textbf{习题43：$p^2q$阶群可解。}$p,q$为不同素数。
\end{problem}
\begin{proof}[补充证明]
由习题25，存在正规西罗子群$P$。
若$|P|=p^2$，由习题20它交换，而$|G/P|=q$，商群为循环群。
若$|P|=q$，$P$循环，而商群阶为$p^2$，仍由习题20交换。

因此两种情况都满足$P$及$G/P$可解。
由可解群的扩张性质，得到$G$可解。
\end{proof}

\begin{problem}
\textbf{习题44：有限$p$群可解。}
\end{problem}
\begin{proof}[补充证明]
对$|G|$归纳。平凡群及$p$阶群都可解。
设$G$为更大的$p$群，类方程说明$Z(G)\neq\{e\}$。
若$Z(G)=G$，$G$交换，当然可解。
否则
\[
1<|G/Z(G)|<|G|.
\]
商群仍为$p$群，归纳假设给出它可解；中心交换，也可解。
由正规子群$Z(G)$及商群的可解性，推出$G$可解。
\end{proof}

\begin{problem}
\textbf{习题45：两个可解商群给出交核的可解商群。}若$H,K\trianglelefteq G$，$G/H,G/K$可解，则$G/(H\cap K)$可解。
\end{problem}
\begin{proof}[补充证明]
定义同态
\[
\varphi:G\to(G/H)\times(G/K),\qquad g\longmapsto(gH,gK).
\]
其核为$H\cap K$，所以第一同构定理给出
\[
G/(H\cap K)\cong\varphi(G)\leq(G/H)\times(G/K).
\]
若两个群$A,B$可解，逐分量计算换位子可得
\[
(A\times B)^{(i)}=A^{(i)}\times B^{(i)}.
\]
取$i$至少为两者导出列终止的步数，右边平凡，所以直积可解。
再用可解群子群可解，得到所求商群可解。
\end{proof}

\subsection*{自同构、生成关系及综合题（习题46--57）}

\begin{problem}
\textbf{习题46：只有两个内自同构时原群交换。}事实上，内自同构群不可能恰有两个元素。
\end{problem}
\begin{proof}[补充证明]
若$|\operatorname{Inn}(G)|=2$，由内自同构群的同构定理，
\[
G/Z(G)\cong\operatorname{Inn}(G)\cong C_2.
\]
商群循环，习题39推出$G$交换。
但交换群的每个内自同构都是单位映射，所以
\[
|\operatorname{Inn}(G)|=1,
\]
与“恰有两个”矛盾。
因此所设条件不可能发生，相应的条件命题成立。
\end{proof}

\begin{problem}
\textbf{习题47：非平凡有限群至少有两个自同构，除了2阶群。}即$|G|>2$时$|\Aut(G)|\geq2$。
\end{problem}
\begin{proof}[补充证明]
\textbf{$G$不交换。}取不在中心的$a$，内自同构$c_a$就不是单位映射，因此至少有两个自同构。

\textbf{$G$交换，且有阶大于2的元素。}取逆映射$x\mapsto x^{-1}$保持乘法，为非单位自同构，结论成立。

\textbf{$G$交换，且每个元素平方为$e$。}
此时以群运算为加法、规定$0x=0,1x=x$，$G$成为$\F_2$上的有限维向量空间。
$|G|>2$说明维数至少2，选一组基$v_1,v_2,\ldots,v_r$。
将$v_1,v_2$交换、其余基向量固定，可唯一线性延拓为非单位的可逆线性映射；它也保持群加法，故为所需自同构。
\end{proof}

\begin{problem}
\textbf{习题48：克莱因四元群的自同构群。}$\Aut(V_4)\cong S_3$。
\end{problem}
\begin{proof}[补充证明]
记$V_4=\{e,a,b,c\}$，其中
\[
a^2=b^2=c^2=e,\qquad ab=ba=c
\]
及其循环替换均成立。
任意自同构固定$e$，因此给出对$\{a,b,c\}$的一个置换。

反过来，任意对这三个元素的置换都保持乘法：相同非单位元相乘得到$e$，不同的两个相乘得到第三个。
把$e$固定后就得到一个群自同构。
复合对应于三点置换的复合，故
\[
\Aut(V_4)\cong S_3.
\]
\end{proof}

\begin{problem}
\textbf{习题49：$S_3,S_4$为完全群。}
\end{problem}
\begin{proof}[补充证明]
由习题6，两群中心都平凡，所以不同元素诱导不同内自同构，分别已有6和24个。

\textbf{$S_3$的情况。}自同构必须置换三个2阶元素。
这三个元素生成$S_3$，所以一个自同构完全由这个置换决定，得到单射
\[
\Aut(S_3)\hookrightarrow S_3.
\]
因而自同构至多6个，内自同构已达到6个，所以全部自同构都是内自同构。

\textbf{$S_4$的情况。}四个西罗3子群分别由固定一个点的3-轮换生成。
它们的正规化子恰为四个点稳定子$H_i\cong S_3$：正规化一个3阶子群须保持它移动的三个点，故固定余下的点；反过来固定该点的全部置换都正规化此3阶子群。
自同构置换这四个正规化子，得到$\Psi:\Aut(S_4)\to S_4$。

若$\alpha\in\ker\Psi$，它固定每个$H_i$。
对任意$g\in S_4$，由$gH_ig^{-1}=H_{g(i)}$，
\begin{align*}
\alpha(g)H_i\alpha(g)^{-1}
&=\alpha(gH_ig^{-1})\\
&=H_{g(i)}.
\end{align*}
左边也等于$H_{\alpha(g)(i)}$，四个稳定子彼此不同，所以$\alpha(g)(i)=g(i)$对每个$i$成立。
故$\alpha(g)=g$，$\Psi$单射。
于是自同构至多24个，而内自同构已有24个，故全部自同构都是内自同构。
结合中心平凡，$S_3,S_4$都是完全群。
\end{proof}

\begin{problem}
\textbf{习题50：非交换单群的自同构群完全。}若$G$为非交换单群，则$\Aut(G)$是完全群。
\end{problem}
\begin{proof}[补充证明]
令$A=\Aut(G)$、$S=\operatorname{Inn}(G)$。
因为$G$非交换且单，中心不能等于$G$，所以$Z(G)=\{e\}$，从而$S\cong G$也单。
由正文无中心群的自同构结论，
\[
C_A(S)=\{e\},\qquad Z(A)=\{e\}.
\]
这里$C_A(S)$表示$A$中与$S$全部元素交换的元素组成的子群。

\textbf{证明$S$在$A$中特征。}
若$\{e\}\neq M\trianglelefteq A$，则$M\cap S\trianglelefteq S$，只能平凡或等于$S$。
若交平凡，由两者在$A$中正规，它们元素的换位子同时落在交中，故$M$与$S$逐元素交换。
于是$M\leq C_A(S)=\{e\}$，矛盾。
因此$S\leq M$，即每个非平凡正规子群都包含$S$。
$S$自身单，所以它是唯一的极小非平凡正规子群；任意自同构必须保持它，故$S$特征。

\textbf{证明$A$的每个自同构是内自同构。}
任取$\beta\in\Aut(A)$，它在$S$上的限制，经$G\cong S$对应于某个$a\in\Aut(G)=A$。
公式$a c_g a^{-1}=c_{a(g)}$说明，$a$在$A$中诱导的内自同构正好与$\beta$在$S$上的限制一致。
将$\beta$与这个内自同构的逆复合，可以假设$\beta$逐元素固定$S$。

对$x\in A,s\in S$，$S\trianglelefteq A$使$xsx^{-1}\in S$，所以
\[
\beta(x)s\beta(x)^{-1}=\beta(xsx^{-1})=xsx^{-1}.
\]
因此$x^{-1}\beta(x)\in C_A(S)=\{e\}$，即$\beta(x)=x$。
所以复合后的自同构恒等，原自同构为内自同构。再结合$Z(A)=\{e\}$，得$A$完全。
\end{proof}

\begin{problem}
\textbf{习题51$^*$：不存在换位子群为$S_4$的更大群。}不存在$G\supsetneq S_4$使$G^{(1)}=S_4$。
\end{problem}
\begin{proof}[补充证明]
反设$N=S_4=G^{(1)}$。$N$为$G$的正规子群，且由习题49为无中心完全群。
任意$g\in G$的共轭限制于$N$为自同构，所以存在$n\in N$使
\[
gxg^{-1}=nxn^{-1}\qquad(\forall x\in N).
\]
这说明$n^{-1}g\in C_G(N)$，因而$G=N C_G(N)$。
两因子逐元素交换，而且
\[
N\cap C_G(N)=Z(N)=\{e\}.
\]
所以乘法映射将$N\times C_G(N)$同构到$G$，从而
\[
G^{(1)}=N^{(1)}\times C_G(N)^{(1)}.
\]
由$G^{(1)}=N$及交平凡，$C_G(N)^{(1)}$必须平凡，故$G^{(1)}=N^{(1)}=A_4$，与$G^{(1)}=S_4$矛盾。
\end{proof}

\begin{problem}
\textbf{习题52：3-轮换为$A_n$中的换位子。}$n>4$时每个3-轮换为一个换位子，从而$A_n^{(1)}=A_n$。
\end{problem}
\begin{proof}[补充证明]
任取三个不同的点$a,b,c$，因为$n\geq5$，还能选不同于它们的$d,e$。
令
\[
x=(ab)(de),\qquad y=(bc)(de).
\]
它们均为偶置换，且$x^{-1}=x,y^{-1}=y$。
直接计算得$xy=(abc)$，因此
\[
[x,y]=x^{-1}y^{-1}xy=(xy)^2=(acb).
\]
交换$x,y$就得到$[y,x]=(abc)$。
所以每个3-轮换都是$A_n$中的换位子。
这些轮换生成$A_n$，故$A_n\leq A_n^{(1)}$；反向包含由定义成立，得到$A_n^{(1)}=A_n$。
\end{proof}

\begin{problem}
\textbf{习题53：$S_4$的两个生成元及关系。}存在$a,b$生成$S_4$，满足$a^2=b^3=(ab)^4=e$。
\end{problem}
\begin{solution}
\textit{补充解答。}
取$a=(12)$、$b=(234)$，按右边先作用计算，有
\[
a^2=e,\qquad b^3=e,\qquad ab=(1234),
\]
所以$(ab)^4=e$。
共轭$a$得到
\[
bab^{-1}=(13),\qquad b^2ab^{-2}=(14).
\]
因此$\langle a,b\rangle$包含$(12),(13),(14)$，它们由习题10生成$S_4$，故$a,b$生成$S_4$。
本题要求构造满足所列关系的生成元，不把这些等式自动当成任意群的完整展示。
\end{solution}

\begin{problem}
\textbf{习题54：$A_4$的两个生成元及关系。}存在$a,b$生成$A_4$，满足$a^2=b^3=(ab)^3=e$。
\end{problem}
\begin{solution}
\textit{补充解答。}
取$a=(12)(34)$、$b=(123)$。
直接计算得
\[
a^2=e,\qquad b^3=e,\qquad ab=(243),
\]
故$(ab)^3=e$。
$a,bab^{-1},b^2ab^{-2}$恰为$V_4$的三个非单位元，所以$\langle a,b\rangle$包含$V_4$。
又$\langle b\rangle$为3阶，与$V_4$交平凡，所以其乘积包含$4\cdot3=12$个元素。
生成群又包含于$A_4$，故只能等于$A_4$。
\end{solution}

\begin{problem}
\textbf{习题55：四元数群的关系与条件。}设$G=\langle a,b\rangle$满足
\[
a^4=b^4=e,\qquad a^2=b^2,\qquad b^{-1}ab=a^{-1}.
\]
若进一步要求$a^2\neq e$，则$G\cong Q_8$；若把这些关系理解为完整展示，也得到$Q_8$。
\end{problem}
\begin{proof}[补充证明]
\textbf{先说明必要限定。}若只要求题中等式成立，$V_4$的两个不同非单位元$a,b$也满足全部关系，却不是8阶群。
因此须补$a^2\neq e$，或把所列关系理解为完整生成关系展示。

\textbf{在$a^2\neq e$下证明阶为8。}
关系$b^{-1}ab=a^{-1}$给出$ab=ba^{-1}$；又$b^2=a^2$，所以任意字可化为
\[
a^ib^j,\qquad0\leq i<4,\quad0\leq j<2.
\]
因此$|G|\leq8$。
$a^2\neq e$及$a^4=e$说明$a$有4阶。
若$b\in\langle a\rangle$，它就与$a$交换，但共轭关系会给$a=a^{-1}$，与$a$有4阶矛盾。
所以$b\notin\langle a\rangle$，两个陪集$\langle a\rangle$、$\langle a\rangle b$不交，各有四个元素，故$|G|=8$。

\textbf{识别为四元数群。}
在$Q_8=\{\pm1,\pm I,\pm J,\pm K\}$中取$I,J$，它们满足相同关系。
将$a^ib^j$对应到$I^iJ^j$，八个正规形两两不同；化简乘积所用关系相同，所以此双射保乘法，给出$G\cong Q_8$。
若采用完整展示读法，展示群至多8阶，同时由$a\mapsto I,b\mapsto J$满射到8阶$Q_8$，故也恰为$Q_8$。

\textbf{列出全部子群并验证正规。}
$Q_8$只有一个2阶元素$-1$，其余六个非单位元均为4阶，分别生成
\[
\langle I\rangle,\qquad\langle J\rangle,\qquad\langle K\rangle.
\]
拉格朗日定理说明真子群阶只能为1、2、4，所以上述三个4阶群连同$\{1\},\{1,-1\}$已列尽所有真子群。
三个4阶子群指数2，均正规；$\{1,-1\}$在中心内，也正规；平凡子群和全群当然正规。
\end{proof}

\begin{problem}
\textbf{习题56$^*$：60以下的群及60阶单群。}每个阶小于60的群都可解；其中单群只可能为素数阶循环群。60阶单群同构于$A_5$。
\end{problem}
\begin{proof}[补充证明]
\textbf{第一步：建立单群的置换作用限制。}
非交换单群$G$的非平凡作用$\rho:G\to S_t$有正规核。
核不能为$G$，所以只能平凡，$\rho$忠实。
符号复合同态$\operatorname{sgn}\circ\rho$若不平凡，同样须单射到$C_2$，与非交换性矛盾。
故像包含于$A_t$，得到
\[
|G|\mid t!/2.
\]
特别地，西罗子群的共轭作用在单群中必须满足$n_p>1$，并要求$|G|\mid n_p!/2$。

\textbf{第二步：排除60以下的非交换单群。}
素数幂阶群有非平凡中心，若单就只能交换，因而只可能为素数阶循环群。
对于其余合数$n<60$，逐项使用$n_p\mid n/p^{v_p(n)}$及$n_p\equiv1\pmod p$。
其中$v_p(n)$表示$n$中素因子$p$的指数。
以下阶数选择其最大素因子$p$就迫使$n_p=1$：
\[
\begin{gathered}
6,10,14,15,18,20,21,22,26,28,33,34,35,38,39,\\
40,42,44,45,46,50,51,52,54,55,57,58.
\end{gathered}
\]
例如$n=45$时$n_5\mid9$，可选1、3、9中只有1模5余1；其余同样由有限个因子核对。
余下仅有$12,24,30,36,48,56$。
12、24、48阶若不已有正规西罗2子群，则$n_2=3$，第一步要求嵌入$A_3$，不可能；36阶若$n_3\neq1$，就有$n_3=4$，要求嵌入$A_4$，也不可能。
30阶单群须$n_5=6,n_3=10$，不同素数阶子群贡献互不重叠的非单位元，至少有
\[
6(5-1)+10(3-1)=44>29,
\]
矛盾。56阶由习题32排除。
所以60以下单群只有素数阶循环群。

\textbf{第三步：证明60以下的全部群可解。}
对阶归纳，平凡群可解。
若$G$单，第二步说明它为素数阶循环群，因而可解。
若不单，取非平凡真正规子群$N$，$N$及$G/N$阶都更小，归纳得二者可解，再由扩张性质得$G$可解。

\textbf{第四步：计算60阶单群的西罗数及元素数。}
设$|G|=60$且$G$单，则由计数及第一步，
\[
n_5=6,\qquad n_3=10.
\]
这里$n_5$只能为1或6；$n_3$只能为1、4、10，1由单性排除，4由$60\nmid|A_4|$排除。
因此5阶元有24个，3阶元有20个，余下非单位元只有
\[
60-1-24-20=15
\]
个。西罗2子群阶为4，其个数只能为1、3、5、15，前两者由单性及$A_3$作用限制排除。

\textbf{第五步：排除$n_2=15$。}
若西罗2子群循环，则15个共轭子群全循环，每个给出两个4阶元；4阶元唯一生成它所属的4阶子群，故不同子群的这些元素不重叠。
于是至少有30个4阶元，超过剩余15个，矛盾。
所以这些西罗子群均为$V_4$，每个含三个对合（即2阶元素），合计45次出现。
全部对合至多15个，所以至少有一个对合$t$属于两个不同的西罗2子群。
这两个$V_4$中的元素都与$t$交换，故它们都包含于$C_G(t)$，使
\[
4\mid|C_G(t)|,\qquad|C_G(t)|>4.
\]
由其阶整除60，只能为12、20或60。
60使$t$属于非交换单群的中心，不可能；20给出指数3的陪集作用，同样由第一步排除。
若为12，则指数5的陪集作用使$G\hookrightarrow A_5$，两群阶相等，故$G\cong A_5$。
但$A_5$的西罗2子群只有五个：其15个对合均为双对换，每个固定唯一一点，固定给定一点的三个双对换加$e$构成一个$V_4$。
对$t=(12)(34)$，$S_5$中与它交换的置换须固定5，可交换两个二元组并在组内交换，所以中心化子阶为$2^2\cdot2!=8$。
其中含奇置换$(12)$，故偶置换部分恰占一半：
\[
|C_{A_5}(t)|=4.
\]
它已经包含固定5的那个$V_4$，所以恰等于它；其他双对换由共轭有相同结论。
若某$V_4$含$t$，就必须包含于这个4阶中心化子，因而等于它。
所以五个这样的子群已列尽，与$n_2=15$矛盾。

\textbf{第六步：用西罗2子群的共轭作用识别群。}
只能$n_2=5$。$G$共轭作用于这五个子群，给出非平凡同态$G\to S_5$。
单性及符号同态的讨论使它单射到$A_5$，再比较阶，得到
\[
G\cong A_5.
\]
\end{proof}

\begin{problem}
\textbf{习题57：有限指数子群仍有限生成。}若$G$有限生成且$[G:H]=m<\infty$，则$H$有限生成。
\end{problem}
\begin{proof}[补充证明]
选有限生成元$g_1,\ldots,g_n$，并把它们的逆元也放入有限集合$S$。
取$H$的$m$个左陪集代表组成$T$，其中$e$代表$H$。
对任意$x\in G$，记$\overline x\in T$为$xH$的代表。

对每个$s\in S,t\in T$，有$stH=\overline{st}H$，因此
\[
b_{s,t}=\overline{st}^{\,-1}st\in H.
\]
这些元素组成有限集合$\mathcal B$，至多$|S|m$个。
写成$a_j^{-1}g_i a_k$时，$a_j$必须选为$g_i a_kH$的代表，不能任取下标$j,k$。

任取$h=s_1\cdots s_r\in H$，$s_i\in S$。
从右端开始令
\[
t_{r+1}=e,\qquad t_i=\overline{s_it_{i+1}}\quad(i=r,r-1,\ldots,1).
\]
归纳可知$t_i$代表后缀$s_i\cdots s_r$所在左陪集；因为$h\in H$，所以$t_1=e$。
令$b_i=t_i^{-1}s_it_{i+1}$，则每个$b_i\in\mathcal B$，且
\begin{align*}
b_1\cdots b_r
&=t_1^{-1}s_1t_2\,t_2^{-1}s_2t_3\cdots t_r^{-1}s_rt_{r+1}\\
&=t_1^{-1}(s_1\cdots s_r)t_{r+1}\\
&=h.
\end{align*}
所以$H$的任意元素都是$\mathcal B$中元素的积，$H=\langle\mathcal B\rangle$有限生成。
\end{proof}
